\section{Introduction}
The rapid scaling of Large Language Models (LLMs) has precipitated an unprecedented increase in the computational and energy demands of artificial intelligence. As the community moves toward massive deployments of generative models, the carbon footprint and operational electricity costs associated with LLM inference have become paramount concerns \cite{schwartz2020greenai, samsi2023words}. To maximize hardware utilization, modern deep learning frameworks increasingly rely on domain-specific languages (DSLs) to generate highly optimized GPU kernels. Among these, OpenAI's Triton \cite{tillet2019triton} has become the \textit{de facto} standard, serving as the compilation backend for PyTorch 2.0, vLLM \cite{kwon2023pagedattention}, and standard libraries like FlashAttention \cite{dao2022flashattention}. 

A core feature of Triton is its native autotuner (\texttt{@triton.autotune}), which sweeps a user-defined grid of hyperparameters---such as \texttt{BLOCK\_SIZE}, \texttt{num\_warps}, and \texttt{num\_stages}---to dynamically select the fastest configuration for a given input tensor shape. However, this optimization process is strictly uni-dimensional: it minimizes execution latency without any consideration for energy consumption. 

Historically, performance engineers have operated under the assumption that "faster is greener"---that the configuration minimizing execution time inherently minimizes energy by returning the GPU to an idle state sooner (the "race-to-sleep" paradigm). While this heuristic holds in some architectural regimes, modern GPUs feature complex thread scheduling, hierarchical memory architectures, and dynamic clock scaling that challenge this assumption. Configurations that maximize throughput often over-provision resources, keeping warps active and memory pipelines saturated in ways that draw disproportionate power relative to their marginal speedup \cite{kernelpro2026, flipflop2026}.

In this work, we present \textbf{GreenTune}, an energy-aware autotuning framework built directly on top of Triton. We instrument the Triton benchmarking loop with high-frequency continuous power sampling to simultaneously measure the latency and energy dissipation of every kernel configuration. By analyzing a representative suite of kernels (Elementwise, Reduction, Normalization, Tiled GEMM, and Fused Causal Attention), we explicitly map the multi-objective Pareto frontiers of modern GPU computations.

This paper makes the following primary contributions:
\begin{itemize}
    \item \textbf{Energy-Aware Autotuning Integration:} We extend the Triton autotuning mechanism to treat energy as a first-class metric, demonstrating that continuous NVML-based power integration can be tightly coupled with microsecond-scale kernel benchmarking.
    \item \textbf{Pareto Characterization of LLM Kernels:} We empirically prove the existence of a non-trivial latency-energy Pareto frontier. We show that memory-bound LLM operations (such as RMSNorm) and compute-bound operations (such as Tiled GEMM) exhibit fundamentally different energy scaling behaviors, often allowing for significant energy savings by selecting a configuration at the Pareto knee rather than the latency extreme.
    \item \textbf{Analysis of Configuration Shape Topology:} We demonstrate that the shape of an energy-optimal configuration systematically differs from a latency-optimal one. Specifically, we observe that energy optimizers aggressively penalize over-partitioned warps (e.g., favoring 16 warps over 32 in RMSNorm) and heavily favor deeper software pipelining (\texttt{num\_stages}: 4 vs. 2 in GEMM) to minimize idle power draw.
    \item \textbf{Cost-Model Guided Search Optimization:} We demonstrate that the immense cost of exhaustive grid searching can be mitigated by utilizing a multi-objective Tree-structured Parzen Estimator (TPE). By leveraging Optuna to explore the search space, we accurately approximate the true Pareto frontier in a fraction of the trials required by exhaustive enumeration.
\end{itemize}

The remainder of this paper is structured as follows. Section \ref{sec:background} provides background on the Triton compiler and GPU power modeling. Section \ref{sec:related_work} contextualizes our work within the broader field of energy-aware compilation. Section \ref{sec:methodology} details the design of GreenTune. Section \ref{sec:setup} describes the experimental environment. Section \ref{sec:results} presents our empirical findings, followed by a discussion of limitations in Section \ref{sec:discussion}, and concluding remarks in Section \ref{sec:conclusion}.
